

% Editorial
\iftrue % change to \iffalse to preview changes
    \newcommand{\saksham}[1]{\textcolor{violet}{\textbf{[Saksham: #1]}}}
    \newcommand{\tianyin}[1]{\textcolor{blue}{\textbf{[ty: #1]}}}
    \newcommand{\rachit}[1]{\textcolor{purple}{\textbf{[rachit: #1]}}}
    \newcommand{\leshna}[1]{\textcolor{pink}{\textbf{[leshna: #1]}}}
    \newcommand{\leshnaedit}[1]{\textcolor{cyan}{#1}}

    \newcommand{\ik}[1]{\textcolor{brown}{\textbf{[ishv1k: #1]}}}
    \newcommand{\ikedit}[1]{\textcolor{magenta}{#1}}
    \newcommand{\ikdelete}[1]{{\color{purple} \sout{#1}}} % del text
    \newcommand{\ikreplace}[2]{{\ikdelete{#1} \ikedit{#2}}} % delete and replace text

    \newcommand{\para}[1]{\smallskip\noindent {\bf #1} }
    
    \newcommand{\todo}[1]{\textcolor{red}{\textbf{[TODO: #1]}}}
    \newcommand{\question}[1]{\textcolor{cyan}{\textbf{[QUESTION: #1]}}}
	\newcommand{\review}[1]{{\color{blue} #1}} % new text

	\newcommand{\JZ}[1]{{\color{orange} #1}} % new text
	\newcommand{\JZx}[1]{{\color{purple} \sout{#1}}} % del text
	\newcommand{\JZc}[1]{{\color{comment} /* \textit{#1} */}} % comment
	\newcommand{\JZd}[1]{} % exile to the void
	\newcommand{\JZe}[1]{{\color{mujica} \uwave{#1}}} % highlight
	\newcommand{\JZhl}{\colorbox{mujica}{\color{white} NOTE:}}
	\newcommand{\JZtodo}{\colorbox{mygo}{\color{white} TODO:}}

	\newcommand{\jz}[1]{{\JZ{#1}}} % new text
	\newcommand{\jza}[2]{{\JZ{#1} \JZc{#2}}} % add and comment text
	\newcommand{\jzc}[2]{{\JZe{#1} \JZc{#2}}} % comment existing text
	\newcommand{\jzd}[2]{{\JZx{#1} \JZc{#2}}} % delete and comment text
	\newcommand{\jzx}[2]{{\JZx{#1} \JZ{#2}}} % delete and replace text
	\newcommand{\jzxc}[3]{{\JZx{#1} \JZ{#2} \JZc{#3}}} % replace text with comment
	\newcommand{\jztodo}[1]{{\JZc{\JZtodo #1}}}
	\newcommand{\jzmark}[1]{{\JZc{\JZhl #1}}}

	\newcommand{\schai}[1]{{\color{blue} #1}} % new text
	% \newcommand{\schai}[1]{#1}
	\newcommand{\schaix}[1]{{\color{red} \sout{#1}}} % del text
	\newcommand{\schaic}[1]{{\color{gray} /* \bf schai \textit{#1} */}} % comment
	
	\newcommand{\jyk}[1]{{\footnotesize{\color{cyan} {\bf JYK: #1}}}}
	\newcommand{\jykx}[1]{{\color{cyan} \sout{#1}}}

	% so we can mix \cite within colored text
	\let\oldcite\cite
	\renewcommand{\cite}[1]{\mbox{\oldcite{#1}}}

\else
	\newcommand{\tianyin}[1]{#1}

    \newcommand{\todo}[1]{}
    \newcommand{\question}[1]{}

	\newcommand{\JZ}[1]{#1} % new text
	\newcommand{\JZx}[1]{} % del text
	\newcommand{\JZc}[1]{} % comment
	\newcommand{\JZd}[1]{} % exile to the void
	\newcommand{\JZe}[1]{} % highlight
	\newcommand{\JZhl}{}
	\newcommand{\JZtodo}{}

	\newcommand{\jz}[1]{#1}
	\newcommand{\jzd}[2]{#2}
	\newcommand{\jzc}[2]{#1}
	\newcommand{\jzx}[2]{#2}
	\newcommand{\jzxc}[3]{#2}
	\newcommand{\jztodo}{}

	\newcommand{\schai}[1]{#1} % new text
	\newcommand{\schaix}[1]{} % del text
	\newcommand{\schaic}[1]{} % comment

	\newcommand{\jyk}[1]{}
	\newcommand{\jykx}[1]{}

	\newcommand{\weiwei}[1]{}
	\newcommand{\weiweix}[1]{}
\fi

\newcommand{\paragraphb}[1]{\vspace{0.1in} \noindent{\textbf{#1}.}}

\newcommand*\circled[1]{\tikz[baseline=(char.base)]{%
            \node[shape=circle,fill=black!20,draw,inner sep=1.5pt] (char) {#1};}}

\newcommand*\circledb[1]{\tikz[baseline=(char.base)]{%
            \node[shape=circle,fill=blue!20,draw,inner sep=1.5pt] (char) {#1};}}